\section{Problem statement}

This preprint addresses exactly one frozen question:

\medskip
\noindent
\textbf{Q9.} \emph{Which enterprise states are decidable under incomplete
observability, and which are provably undecidable, in enterprises modeled
as continua $K_{EA}$?}
\medskip

The question is posed in the formal setting of the Ontology of Continua
(OC Core) and the executable specification ETS.K\_EA.v1.1. It is treated
as a problem in logical decidability, not as a problem of estimation,
reconstruction, inference, or control.

\subsection{Scope of the question}

An \emph{enterprise state} is understood here not as a full internal
configuration of the enterprise continuum, but as a state predicate
defined over admissible OC/ETS elements. In particular, this includes:
\begin{itemize}
\item admissible phase outcomes (OK, INERTIA, COLLAPSE, STOP);
\item predicates defined via thresholds $\Theta_{EA}$, boundaries
      $\partial\Omega_{EA}$, or structural properties of cycles and
      flows;
\item any property expressible without introducing new primitives
      beyond those fixed in OC Core and ETS.K\_EA.v1.1.
\end{itemize}

No commitment is made to any operational, managerial, or interpretative
meaning of such states beyond their formal definition.

\subsection{Incomplete observability}

Enterprises are assumed to be only partially observable. Observations
are restricted to the admissible observability constructs already
defined in ETS, yielding observable signatures in
$\Sigma_{\mathrm{obs}}$ via the observation operator $\Pi$.
No assumption is made that observable information is sufficient to
reconstruct internal states, nor that it becomes sufficient with more
data or time.

Incomplete observability is therefore taken as a structural condition,
not as a temporary limitation or a noise-related issue.

\subsection{Decidability versus observability}

A central constraint of this work is the strict separation between
observability and decidability. Observability concerns what can be
measured or registered as an observable signature. Decidability
concerns whether a state predicate has a truth value that is invariant
across all enterprise continua consistent with the same observable
signature.

Accordingly:
\begin{itemize}
\item high observability does not imply decidability;
\item undecidability is not a defect of methods, data quality, or
      inference techniques;
\item decidability and undecidability are treated as intrinsic logical
      properties induced by the OC/ETS structure.
\end{itemize}

\subsection{Negative framing}

The problem is deliberately framed to admit negative results. The
objective is not to maximize the set of decidable states, but to
classify, within the fixed formalism, which classes of predicates are:
\begin{itemize}
\item decidable under incomplete observability;
\item undecidable under incomplete observability;
\item decidable only in a STOP-only sense, at the diagnosability
      boundary.
\end{itemize}

No prescriptions, governance recommendations, optimization criteria,
or intervention strategies are derived from these results. The outcome
is a classification of logical limits, not a method for overcoming
them.
